\begin{ouverture}[statistiques]{7}{Proportions et pourcentages}{Partie 2 · Modéliser et mesurer}
{Centièmes · Proportions · Appliquer un pourcentage}
\objectif{Donner du sens à une écriture en pourcentage.}
\objectif{Exprimer une proportion simple sous forme de fraction, de décimal ou de pourcentage.}
\objectif{Appliquer un pourcentage à un nombre ou à une grandeur.}
\prerequis{Fractions décimales ; fraction d’un nombre ; multiplication par un entier.}
\end{ouverture}
\pagestyle{otMain}\essentiel\fontsize{11}{15}\selectfont
\phantomsection\label{cours-7-1}
\essSection{1}{La référence à cent}
\begin{definition}Un pourcentage est une fraction de dénominateur 100 : 35~\% $=\frac{35}{100}=0{,}35$. « 35 pour cent » signifie 35 sur 100 : 35 parts d’un tout partagé en 100 parts égales.\end{definition}
\label{fig-cours-7-1-cours-grille-100}\figureGeo{cours-grille-100}

\essSection{2}{Relier fraction, décimal et pourcentage}
\begin{propriete}Pourcentages à connaître : 100~\% représente le tout ; 50~\% $=\frac12=0{,}5$ ; 25~\% $=\frac14=0{,}25$ ; 75~\% $=\frac34=0{,}75$ ; 10~\% $=\frac1{10}=0{,}1$ ; 20~\% $=\frac15=0{,}2$.\end{propriete}
\begin{savoirfaire}[Changer d’écriture]Pour écrire $\frac35$ en pourcentage, on cherche une fraction égale de dénominateur 100 : $\frac35=\frac{60}{100}=60$~\%. Inversement, 8~\% $=\frac8{100}=0{,}08$.\end{savoirfaire}
\clearpage
\phantomsection\label{cours-7-2}
\essSection{3}{Exprimer une proportion simple}
\begin{definition}La proportion d’une partie dans un tout est le quotient de la partie par le tout. Sur 25 élèves, 9 portent des lunettes : la proportion est $\frac9{25}=\frac{36}{100}$, soit 36~\%.\end{definition}
\begin{savoirfaire}[Se ramener à cent]Quand le tout est un diviseur de 100 (2, 4, 5, 10, 20, 25, 50) ou un multiple de 100, on obtient facilement un dénominateur 100 : $\frac7{20}=\frac{35}{100}=35$~\% ; $\frac{12}{200}=\frac6{100}=6$~\%.\end{savoirfaire}
\essSection{4}{Appliquer un pourcentage}
\begin{savoirfaire}[Calculer un pourcentage d’un nombre]Prendre 30~\% de 80, c’est calculer $\frac{30}{100}\times80=24$. On peut aussi passer par 10~\% : 10~\% de 80 vaut 8, donc 30~\% de 80 vaut $3\times8=24$.\end{savoirfaire}
\label{fig-cours-7-2-cours-bande-pct}\figureGeo{cours-bande-pct}

\begin{attention}Un pourcentage se rapporte toujours à un tout : 50~\% de 20 € et 50~\% de 200 € ne sont pas la même somme. Une partie d’un tout représente au plus 100~\% de ce tout.\end{attention}

\clearpage\rubrique{Découverte}
\begin{blocLivre}{ex-7-1}\begin{exercice}[1]{Lire une grille}
La grille compte 100 cases. Quelle proportion de la grille est colorée ? L’écrire sous forme de fraction, d’écriture décimale et de pourcentage.
\label{fig-ex-7-1}\figureGeo{pct-1}
\end{exercice}\end{blocLivre}
\begin{blocLivre}{ex-7-2}\begin{exercice}[1]{Trois écritures}
Recopier et compléter le tableau.\begin{center}\renewcommand{\arraystretch}{1.35}\begin{tabular}{c|cc}\toprule
Fraction & Écriture décimale & Pourcentage\\\midrule
$\frac34$ & … & …\\
… & $0{,}1$ & …\\
… & … & 60 \%\\
$\frac15$ & … & …\\
\bottomrule\end{tabular}\end{center}
\end{exercice}\end{blocLivre}
\begin{blocLivre}{ex-7-3}\begin{exercice}[1]{Pourcentages usuels}
Calculer de tête : 50~\% de 46 ; 25~\% de 80 ; 10~\% de 350 ; 75~\% de 12.
\end{exercice}\end{blocLivre}
\begin{blocLivre}{ex-7-4}\begin{exercice}[1]{Une proportion}
Dans une classe de 25 élèves, 9 portent des lunettes. Quelle est la proportion d’élèves qui portent des lunettes ? L’écrire en pourcentage.
\end{exercice}\end{blocLivre}
\begin{blocLivre}{ex-7-5}\begin{exercice}[1]{Partie d’un tout}
Un sac contient 200 billes ; 30~\% d’entre elles sont rouges. Combien y a-t-il de billes rouges ?
\end{exercice}\end{blocLivre}

\rubrique{Entraînement}
\begin{blocLivre}{ex-7-6}\begin{exercice}[2]{Des proportions simples}
Exprimer en pourcentage : 7 sur 20 ; 12 sur 50 ; 3 sur 4 ; 45 sur 300.
\end{exercice}\end{blocLivre}
\begin{blocLivre}{ex-7-7}\begin{exercice}[2]{Appliquer à une grandeur}
Calculer : 20~\% de $3{,}5$ kg ; 15~\% de 60 min ; 5~\% de 240 €.
\end{exercice}\end{blocLivre}
\begin{blocLivre}{ex-7-8}\begin{exercice}[2]{Les trajets du matin}
La bande représente les 240 élèves d’un collège selon leur moyen de transport. Quel pourcentage vient à pied ? Combien d’élèves cela représente-t-il dans chaque catégorie ?
\label{fig-ex-7-8}\figureGeo{pct-8}
\end{exercice}\end{blocLivre}
\begin{blocLivre}{ex-7-9}\begin{exercice}[2]{Comparer des proportions}
Dans la classe A, il y a 18 filles sur 30 élèves ; dans la classe B, 13 filles sur 25 élèves. Dans quelle classe la proportion de filles est-elle la plus grande ?
\end{exercice}\end{blocLivre}
\begin{blocLivre}{ex-7-10}\begin{exercice}[2]{Le pourcentage manquant}
Lors d’un sondage auprès de 500 personnes, 46~\% ont répondu oui et 38~\% non ; les autres sont sans avis. Quel pourcentage est sans avis ? Combien de personnes cela fait-il ?
\end{exercice}\end{blocLivre}

\rubrique{Problèmes}
\begin{blocLivre}{ex-7-11}\begin{exercice}[3]{Les soldes}
Un jean coûte 60 €. Pendant les soldes, son prix baisse de 30~\%. Calculer la réduction, puis le nouveau prix.
\end{exercice}\end{blocLivre}
\begin{blocLivre}{ex-7-12}\begin{exercice}[3]{Le jus de fruits}
Une brique de $1{,}5$ L de boisson contient 12~\% de jus de pomme. Quel volume de jus de pomme contient-elle ? Répondre en litres, puis en centilitres.
\end{exercice}\end{blocLivre}
\begin{blocLivre}{ex-7-13}\begin{exercice}[3]{La commune et sa forêt}
Une commune s’étend sur $2\,400$ hectares : 35~\% sont couverts de forêt, 40~\% de champs, et le reste par des habitations et des routes. Calculer la surface de chaque partie.
\end{exercice}\end{blocLivre}
\begin{blocLivre}{ex-7-14}\begin{exercice}[3]{L’élection des délégués}
Dans une classe, 28 élèves ont voté : Inès a obtenu 14 voix, Jules 7 voix et Sami 7 voix. Exprimer le résultat de chaque candidat en pourcentage.
\end{exercice}\end{blocLivre}
\begin{blocLivre}{ex-7-15}\begin{exercice}[3]{Vrai ou faux}
Maël affirme : « 20~\% de 50, c’est la même chose que 50~\% de 20. » Vérifier par le calcul, puis expliquer.
\end{exercice}\end{blocLivre}

\clearpage
